%% frag_benchmark.tex
%% Insertion point: Section 3, immediately before \subsection{Experimental Setup}.

\subsection{A Multi-Ecosystem Benchmark}\label{sec:benchmark}

The evaluation in the conference version used 50 questions drawn from a single
community source, which is too few to separate an architectural effect from
noise and too narrow to say anything about ecosystems other than the ones those
questions happened to mention. We therefore constructed a larger benchmark, and
we describe its construction here in enough detail to expose its biases, because
a benchmark assembled by the authors of the system it evaluates warrants that
scrutiny.

\subsubsection{Construction}

The benchmark contains 300 questions arranged as a $24 \times 5$ grid of
software ecosystem by question category, filled to 12 or 13 questions per
ecosystem and exactly 60 per category. The 24 ecosystems span mobile and desktop
platforms (Apple iOS, Apple macOS, Android, Windows), four Linux distributions
(Ubuntu, Debian, Fedora, Arch) and the kernel itself, browsers (Firefox,
Chrome), container infrastructure (Docker, Kubernetes), language package
ecosystems (npm/Node, pip/Python, cargo/Rust), self-hosted and home applications
(Home Assistant, NAS/self-hosted, WordPress), databases (PostgreSQL, MySQL,
NoSQL), developer tooling, and LLM/AI model releases. The five categories are
releases, bugs, security, community, and general.

Every question is mined from live data rather than authored: 249 are verbatim
titles of real community posts, and 51 are generated from real release or
advisory records. No question is a template with a hand-written stem. The
builder draws from a pool of 9,923 community posts and the release and advisory
endpoints described in Section~\ref{sec:datasources}, and it is deterministic:
a fixed seed, an on-disk response cache, and an \texttt{--offline} mode that
replays the cache, so two runs of the builder produce byte-identical output.
Each record carries the endpoint and record type it came from, so mined and
generated questions remain distinguishable in any analysis.

\subsubsection{Two Upstream Data Defects}

Two properties of the upstream feeds would have injected false questions into
the benchmark had they been consumed naively. We report them because they are
properties of the public data sources rather than of our pipeline, and anyone
building on the same feeds will encounter them.

First, on advisory records the product field names the \emph{index token} under
which the advisory was filed, not the affected product. An advisory for ``Mocha
Telnet Lite for iOS 4.2'' is stored with product \texttt{iOS} and version
\texttt{4.2.0}; advisories concerning Velero are stored under
\texttt{Kubernetes}. Templating those fields directly yields questions such as
\emph{``Is iOS v4.2.0 vulnerable?''}, which assert a falsehood in the question
itself and would be scored against a system for failing to confirm it. The
builder instead parses the affected product and version out of the advisory
text, keeps a record only when the parsed product belongs to a tracked
ecosystem, and discards the rest: of 3,113 advisory rows, 548 parsed and 268
could be confidently attributed.

Second, a minority of records carry malformed dates (for example
\texttt{23.30.13102017} and \texttt{2025[29]0213}). Since date agreement is one
of the two decisive facts in this domain, these records are dropped rather than
normalized by guessing.

We also stopped trusting the feed's own topical flags. Its CVE-relatedness flag
fires on posts such as \emph{``Best budget AI subscription / API tokens?''}, so
it is used only to break ties within the security category rather than to assign
one.

\subsubsection{Ground Truth and Its Limits}

Of the 300 questions, 51 carry a reference answer, and only those generated from
release or advisory records do --- a community post title states a question but
does not come with an authoritative answer. Ground-truth coverage is therefore
concentrated in the releases and security categories, and any correctness or
hallucination measurement over this benchmark is computed on that subset rather
than on all 300 questions.

This has a consequence for sampling that we did not anticipate and that a
reader reusing the benchmark should know. A subset selected to balance ecosystem
and category, without regard to ground truth, can leave almost none: our first
100-question stratified selection carried a reference answer on 2 questions. A
selection that prefers ground-truth records within each grid cell retains 28.
Reported retrieval metrics are unaffected --- they rest on pooled relevance
judgments, not reference answers --- but any answer-correctness figure requires
the second kind of selection, and we report retrieval and answer-faithfulness
results rather than correctness for this reason.

\subsubsection{Coverage Gaps}

One of the 120 grid cells (Arch $\times$ security) had no live candidate
available at build time. Rather than synthesize one, the builder logs the gap
and fills the category quota from other ecosystems, so the category totals hold
exactly while that single cell is empty. Nineteen product-name lookups returned
404 because the product is not in the vendor catalogue (npm, pip, cargo,
Synology, Home Assistant, NixOS, GitLab among them); this is a catalogue
coverage limitation rather than an endpoint failure, and the affected ecosystems
are represented through their community and release records instead.
